\section{Implementing the field encodings}\label{sec:encoders}

A matrix with few nonzero entries in each row admits a useful slot
construction. First prepare each column while retaining its input index.
Then, conditional on the output row, erase that index with an amplitude
determined only by the maximum row size. Two reflections amplify the
result. When the matrix is nearly isometric, all its singular directions
require almost the same amplification angle.
The reflection argument is amplitude amplification
\cite[Section~2]{BHMT2002}. Preparing states to realize a matrix as a block
of a larger unitary is also a standard block-encoding method
\cite[Lemmas~47--48]{GSLW2019}. Here arbitrary table-controlled column
preparation permits us to use sparsity only in the rows.

\subsection{A construction for a nearly isometric matrix}

Let $\mathcal U$ be a set of rows encoded in a $w$-qubit register, and
let $T:\C^d\to\C^{\mathcal U}$ have entries $a_{ui}$, where $d\le2^n$.
Fix an integer $L\ge1$. Suppose that every row has at most $L$ nonzero
entries and
\begin{equation}\label{encode:assumptions}
 (1-\varepsilon)I_d\le T^*T\le I_d,
 \qquad L\ge1,\quad 0\le\varepsilon<1.
\end{equation}
All these data may depend on a label register, with the same dimensions
and bounds for every label. The construction below treats labels as controls.

Extend the $n$-qubit input index by
$h=\lceil\log_2(L+1)\rceil$ high bits, and add the row register and two
one-bit flags. Let $\iota$ embed $\C^d$ with all these new qubits zero.
Let $\jmath$ embed $\C^{\mathcal U}$ with the extended index zero and
both flags equal to one. Write $P_0=\iota\iota^*$ and
$Q=\jmath\jmath^*$ for the corresponding projections. Thus $Q$ allows
any valid row, while $P_0$ allows any of the $d$ input indices.
We call the range of $Q$ the \emph{good subspace}.

\begin{lemma}\label{encode:row-sparse}
Under \eqref{encode:assumptions}, an exact-slot unitary $\mathsf U_T$
satisfies
\begin{equation}\label{encode:generic-error}
 \norm{\mathsf U_T\iota-\jmath T}\le5\varepsilon.
\end{equation}
It uses $O(L(n+w+h))$ slots and $w+h+2$ fresh qubits. Its slot structure
depends only on the register sizes, $d$, and $L$, and it works coherently
for any family of matrices indexed by labels.
\end{lemma}

\begin{proof}
Set
\begin{equation}\label{encode:angle}
 \theta=\frac\pi{4L+2},\qquad t_0=\sin\theta.
\end{equation}
We first construct a unitary $\mathsf B$ with the prescribed block
\begin{equation}\label{encode:block}
 \jmath^*\mathsf B\iota=t_0T.
\end{equation}
For each valid input $i$, the column norm is at most one. Conditional on
$i$, use Lemma~\ref{slot:preparation} to prepare the row register and
first flag in a unit vector whose flag-one part is
\[
 \sum_{u\in\mathcal U}a_{ui}\ket u\ket1.
\]
The remaining squared norm is placed in the flag-zero subspace. The
input index is unchanged during this preparation.

For each row $u$, choose a list $I_u$ of exactly $L$ distinct extended
indices containing every $i$ with $a_{ui}\ne0$. Such lists exist because
the extended index space has dimension
$2^{n+h}\ge2^n(L+1)>L$. Padding indices need not be valid inputs.
Conditional on the row, apply the inverse of a state preparation for
\[
 \frac1{\sqrt L}\sum_{i\in I_u}\ket i
\]
on the extended index register. Its matrix element from any $i\in I_u$
to zero is $1/\sqrt L$. Finally rotate the second flag so that its
amplitude at one is $t_0\sqrt L$. This is possible because
$t_0\sqrt L\le\theta\sqrt L\le1$. Prescribe arbitrary unitary extensions
for invalid indices and rows. Projecting onto index zero and flags one
now multiplies each relevant $a_{ui}$ by
$(1/\sqrt L)(t_0\sqrt L)=t_0$, proving \eqref{encode:block}.
All three operations together form $\mathsf B$.

Starting with $\mathsf B$, perform $L$ rounds of
\begin{equation}\label{encode:round}
 (2\mathsf B P_0\mathsf B^*-I)(I-2Q).
\end{equation}
These are implementable with $\mathsf B$, its inverse, and diagonal
reflections. We verify their action on every singular direction of $T$.
Let $v_j$ be an orthonormal basis of right singular vectors and write
\[
 Tv_j=\sigma_j e_j,\qquad
 \sqrt{1-\varepsilon}\le\sigma_j\le1,
\]
where the $e_j$ are orthonormal. In this calculation identify the input
and output spaces with their embeddings $\iota$ and $\jmath$.
By \eqref{encode:block},
\begin{equation}\label{encode:plane}
 \mathsf Bv_j=\sin\beta_j\,e_j+\cos\beta_j\,e'_j,
 \qquad \beta_j=\arcsin(t_0\sigma_j),
\end{equation}
for a unit vector $e'_j$ in the kernel of $Q$. The $e'_j$ are mutually
orthogonal: this follows by taking pairwise inner products in
\eqref{encode:plane}, using unitarity of $\mathsf B$ and orthogonality of
the $e_j$. Each $e'_j$ is also orthogonal to every $e_i$.

For completeness, both reflections preserve the plane spanned by
$e_j,e'_j$. The first reflection is about the range of $\mathsf B P_0$,
which is spanned by the mutually orthogonal vectors $\mathsf Bv_i$.
Within the $j$th plane it is therefore reflection about
$\sin\beta_j\,e_j+\cos\beta_j\,e'_j$; the second reflection changes the
sign of $e_j$ and fixes $e'_j$. Their product in the order
\eqref{encode:round} sends
\[
 \sin\alpha\,e_j+\cos\alpha\,e'_j
 \quad\hbox{to}\quad
 \sin(\alpha+2\beta_j)e_j+\cos(\alpha+2\beta_j)e'_j.
\]
After $L$ rounds, the angle is $(2L+1)\beta_j$. As
$(2L+1)\theta=\pi/2$, the distance of this vector from $\sigma_j e_j$
is at most
\begin{align*}
 (2L+1)(\theta-\beta_j)+(1-\sigma_j)
 &\le (2L+1)\,2t_0(1-\sigma_j)+(1-\sigma_j)\\
 &\le(\pi+1)(1-\sigma_j)
 \le5\varepsilon.
\end{align*}
Here $t_0\le1/2$, so the derivative of $\arcsin$ is at most $2$ on
$[0,t_0]$, and
$1-\sigma_j\le1-\sqrt{1-\varepsilon}\le\varepsilon$.
The error vectors lie in mutually orthogonal planes. The same bound
therefore holds in operator norm on the whole input space, proving
\eqref{encode:generic-error}.

The first preparation uses $O(w+1)$ slots and the second $O(n+h)$, so
$\mathsf B$ uses $O(n+w+h)$ slots. The complete circuit contains
$2L+1$ occurrences of $\mathsf B$ or $\mathsf B^*$ and $2L$ diagonal
reflections. A diagonal reflection is one slot, as explained in
Section~\ref{sec:slots}. No additional workspace is needed. All prescriptions
can be made separately for each label using the same slot positions and
controls. Since the resulting maps are block diagonal in the labels,
the bounds hold for coherent superpositions of labels as well.
\end{proof}

\subsection{Applying the construction to the Fourier fields}

In the preceding section, a row is a frequency and a port. Let
\begin{equation}\label{encode:register-size}
 K=b(L+1),\qquad h=\lceil\log_2(L+1)\rceil.
\end{equation}
Encode the frequency in $K$ qubits and the port in $n$ qubits, using
the first $k$ port states. The two field maps have the same row space
and therefore use the same output embedding $\jmath$: extended input
index zero, both flags one, and a valid port. This shared layout will
allow the second encoding to be reversed after the recursive call.

\begin{proposition}\label{encode:implementation}
Under the hypotheses and definitions of Section~\ref{sec:fields}, there
are exact-slot unitaries $\mathsf U_E,\mathsf U_R$ such that
\begin{equation}\label{encode:field-error}
 \norm{\mathsf U_E\iota-\jmath\mathcal E_L}\le5r^{2L},
 \qquad
 \norm{\mathsf U_R\iota-\jmath\mathcal R_L}\le5r^{2L}.
\end{equation}
Each uses $O\bigl(L(K+n+h)+K\bigr)$ slots and $K+n+h+2$ fresh qubits,
with slot structure independent of $V$. The construction works coherently
with any additional label register.
\end{proposition}

\begin{proof}
Write either field map in Fourier coordinates. Proposition~\ref{fields:sparsity}
gives at most $L$ nonzero entries in each row, and
Proposition~\ref{fields:gram} supplies \eqref{encode:assumptions} with
$\varepsilon=r^{2L}$. Apply Lemma~\ref{encode:row-sparse}, taking
$w=K+n$, to obtain the corresponding coefficient map with error at most
$5r^{2L}$. Then apply the positive Fourier transform to each of the
$L+1$ frequency factors, so that $\ket w$ becomes
$|\mathcal G|^{-1/2}\sum_z\chi_w(z)\ket z$.
The same physical register now holds the phase label $z$, rather than the
Fourier frequency $w$. Lemma~\ref{slot:fourier} implements these transforms
using $O(K)$ slots.
They preserve the range of $Q$, because every frequency is allowed there
and the index, flags, and port are unchanged. Thus they convert the
coefficient map into the field map in the same embedding $\jmath$,
while preserving the error norm on the full space. Their addition
gives \eqref{encode:field-error} and the stated resource bounds.
\end{proof}
